\documentclass[11pt,a4paper]{article}
\usepackage[T1]{fontenc}
\usepackage{lmodern,amsmath,amssymb,amsthm,mathtools,booktabs,microtype}
\usepackage[margin=27mm]{geometry}
\usepackage[colorlinks=true,linkcolor=blue,citecolor=blue,urlcolor=blue]{hyperref}
\setlength{\emergencystretch}{2em}
\newtheorem{theorem}{Theorem}[section]
\newtheorem{lemma}[theorem]{Lemma}
\newtheorem{corollary}[theorem]{Corollary}
\newcommand{\Q}{\mathbb Q}\newcommand{\Z}{\mathbb Z}
\DeclareMathOperator{\N}{N}
\title{Actual detector neighborhoods without square complements}
\author{}\date{16 November 2026}
\begin{document}\maketitle
\begin{abstract}
For every odd coefficient exponent in the native general detector source,
one of the integer inputs $3,6,7$ has complete guarded prepared witness
families in which every pole complement $17(17-b_j)$ is a nonsquare in
$\Q_{23}$. The obstruction persists throughout Cartesian families of
independent auxiliary-index returns. A uniform eleven-row residue
certificate proves the assertion. It also applies to the explicit
three-detector source. Consequently shrinking arbitrary prepared return
neighborhoods cannot impose rational square complements. Other actual
witness families and common rational pole fibers are not excluded.
\end{abstract}

\section{The actual scalar and detector functions}

The native general source \cite{General,UnitNorm} uses an odd fixed
coefficient exponent $L\ge3$, the prime $\ell=76249$, and
\begin{align}
 \beta&=216X^4+450459,\label{eq:beta}\\
 r&=\beta(\beta-1)(\beta-2401)(32\beta+343)(2\beta-1),\label{eq:r}\\
 \rho&=\frac{102^2\ell^Lr}
       {[1+17(\ell^{(L+1)/2}+1)r]^2}.\label{eq:rho}
\end{align}
The actual detector values are
\begin{equation}\label{eq:detector}
 b_j=(\rho R_j)
       \left(\frac{\N_{K/\Q}(W_j)}{(W_j)_0}\right)^2,
 \qquad K=\Q(\sqrt2).
\end{equation}
There are two detectors in the general twenty-five-variable source and
three in the explicit twenty-seven-variable source \cite{Explicit}.
Their contact functions and comparison scales differ, but the scalar
\eqref{eq:rho} is the same. The notation $(W_j)_0$ denotes the rational
component of the quadratic-field value $W_j$.

We retain the arithmetic premises and precise completeness conclusions
of these sources. A prepared actual input has square reduced denominator,
lies in $(0,17)$, is a square in $\Q_2,\Q_3,\Q_{17}$, and has the
stated minimum valuations $3,1,2$ there. Point-stage guards are required
to be nonzero at every geometric character of the finite point algebra.

\begin{theorem}\label{thm:actual}
Under the retained source premises, for every odd $L\ge3$ there is an
integer $X\in\{3,6,7\}$ and a guarded prepared actual point-stage tuple
of the general source such that
\begin{equation}\label{eq:nonsquare}
                  17(17-b_j)\notin\Q_{23}^{\times2}
                  \qquad(j=0,1).
\end{equation}
The fixed input, root points, comparison slope, and scalar data can be
retained while the five auxiliary cyclic indices vary independently in
infinite return sets. Every resulting tuple remains guarded and prepared
and still satisfies~\eqref{eq:nonsquare}. Each can be completed by its
ordinary individual pole witnesses.

The same statement holds for all three detector values in the explicit
source. At its certified exponent $L=20483$, one may take $X=6$; all
three complements then equal $7$ modulo $23$.
\end{theorem}

The theorem concerns a complete prepared family for each indicated
integer input. It does not assert that every actual witness for that
input belongs to this family. In particular it is not a counterexample
to the existence of some other actual tuple with square complements.

\section{A uniform residue certificate}

Put $k=(L+1)/2$. Modulo $23$ one has
\[
 \ell=4,\qquad \operatorname{ord}_{23}(4)=11,\qquad
 102^2=8,\qquad \beta=9X^4+4.
\]
Thus the scalar residues depend only on $k$ modulo $11$ and $X$ modulo
$23$. Write $d=1+17(4^k+1)r$ for the denominator before squaring.

\begin{lemma}\label{lem:table}
Every row in the following table has $r,d,\rho$ nonzero modulo $23$,
and its final entry is a quadratic nonresidue modulo $23$.
\begin{center}
\normalfont
\begin{tabular}{rrrrrrr}\toprule
$k\bmod11$&$X$&$\beta$&$r$&$d$&$\rho$&$17(17-\rho)$\\\midrule
0&3&20&1&12&8&15\\
1&6&7&19&6&22&7\\
2&6&7&19&18&11&10\\
3&3&20&1&2&1&19\\
4&7&16&9&15&9&21\\
5&3&20&1&15&16&17\\
6&3&20&1&6&13&22\\
7&6&7&19&10&5&20\\
8&7&16&9&13&4&14\\
9&3&20&1&9&9&21\\
10&3&20&1&5&13&22\\\bottomrule
\end{tabular}
\end{center}
\end{lemma}
\begin{proof}
Reduce~\eqref{eq:beta}--\eqref{eq:rho} and substitute the displayed
input. The scalar residue is
\[
                \rho=8\,4^{2k-1}r/d^2\pmod{23}.
\]
Direct exponentiation gives $4^{11}=1$ and $4^j\ne1$ for
$1\le j<11$. The nonzero square residues modulo $23$ are
\[
                    1,2,3,4,6,8,9,12,13,16,18.
\]
No final-column value belongs to this list. The companion exact checker
also evaluates the native scalar coefficient circuits for representatives
of all eleven exponent classes.
\end{proof}

\section{Compatibility with complete actual source witnesses}

\begin{proof}[Proof of Theorem~\ref{thm:actual}]
Fix the exponent and choose the input from Lemma~\ref{lem:table}.
The prime $23$ belongs to the original fixed preparation set: the general
source includes all primes at most $76176$ before choosing roots, the
slope, or the fresh comparison primes. Its scalar preparation and slope
selection permit $R_j$ and both conjugate values of $W_j$ to be as close
to one as prescribed at every such prime
\cite[Section 5]{General}. Impose, in particular,
\begin{equation}\label{eq:congruences}
 R_j\equiv1\pmod{23},\qquad
 W_{j,+}\equiv W_{j,-}\equiv1\pmod{23}.
\end{equation}
The prime splits in $K$, since $5^2=2\pmod{23}$; the two displayed
values may be read in its two completions. Both their product and their
half-sum are one modulo $23$. Formula~\eqref{eq:detector} therefore gives
\begin{equation}\label{eq:rho-residue}
                           b_j\equiv\rho\pmod{23}.
\end{equation}
Lemma~\ref{lem:table} now shows that every $17(17-b_j)$ is a nonsquare
unit in $\Q_{23}$.

These congruences fit the original order of choices. Prepare the scalar
functions and the required finite precision, choose root points and
comparison scales, then apply the prime-pattern theorem with sufficiently
deep denominator divisibility at $23$, and finally choose the auxiliary
index in the prescribed compatible progression. No height scale is made
to depend on a point chosen later. The usual argument makes the resulting
actual detector inputs prepared at every prime and at the real place.

The full-character guard proposition in \cite[Section 5]{UnitNorm}
preserves all previously prepared finite-place conditions while removing
the finitely many real guard zero sets. Hence it preserves
\eqref{eq:congruences} and \eqref{eq:rho-residue}. This supplies a guarded
prepared actual point-stage tuple with the required local nonsquares.

Apply the independent-index theorem of \cite[Theorem 3.1]{Returns}, adding
\eqref{eq:rho-residue} to its prescribed finite congruences. It gives a
positive integer $M$ and an infinite set $\mathcal T\subset\Z_{>0}$ such
that independent auxiliary choices
\[
                   R_i=[m_{i,0}+2Mt_i]P,
                   \qquad(t_1,\ldots,t_5)\in\mathcal T^5,
\]
remain guarded and prepared with all the same fixed earlier data. Their
complements stay nonsquare in $\Q_{23}$. The prepared-input pole
realization supplies separate ordinary pole witnesses for each tuple;
the established pole-density argument supplies their native guards.

The explicit source uses the same scalar preparation and has three
contact functions, all tending to one at the preparation places
\cite[Section 6]{Explicit}. The same argument applies to all three.
Finally $(20483+1)/2\equiv1\pmod{11}$, and the second table row gives
$X=6$, $\rho\equiv22$, and $17(17-b_j)\equiv7\pmod{23}$.
\end{proof}

\begin{corollary}
The independent-index preparation theorem cannot be strengthened to
assert that every prepared return neighborhood contains an actual tuple
with rational square complements. Shrinking the neighborhoods constructed
in Theorem~\ref{thm:actual} cannot produce such a tuple.
\end{corollary}
\begin{proof}
A nonzero rational square is a square in $\Q_{23}$. The required
congruence neighborhood contains only nonsquare complements there.
\end{proof}

This obstruction concerns the square-complement equations. It does not
obstruct points on the original pole curves: local solubility at every
rational shared parameter holds for every prepared input
\cite{Local}. Nor does it exclude a common rational fiber obtained from
the same or a different prepared actual tuple.

\section{Exact source binding and verification}

The verifier is
\begin{center}\small
\texttt{negative\_results/computation/}\\
\texttt{actual\_square\_complement\_v43/check.py}.
\end{center}
It checks the eleven rows with exact rational arithmetic and independent
modular exponentiation, traverses eleven native general scalar circuits,
and separately checks the explicit source at $L=20483$. It records the
hashes of the coefficient builder and the supporting source papers. The
infinite existence of the actual prepared families uses precisely the
retained source premises and the cited completeness and return theorems.

\begin{thebibliography}{9}
\raggedright
\bibitem{General}
\emph{A defining polynomial in twenty-five variables: combined scalar
and comparison conditions}, project source, version 19,
\texttt{papers/few\_variables25/}\allowbreak
\texttt{few\_variables25\_merged\_20261110\_v19.tex}.
\bibitem{Explicit}
\emph{An explicit defining polynomial in twenty-seven variables},
project source, version 19,
\texttt{papers/few\_variables27/}\allowbreak
\texttt{explicit27\_20261110\_v19.tex}.
\bibitem{UnitNorm}
\emph{A single norm determinant for the 25- and 13-variable constructions},
project source, version 28,
\texttt{papers/finite\_algebra/}\allowbreak
\texttt{unit\_norm\_20261114\_v28.tex}.
\bibitem{Returns}
\emph{Fixed pole fibers and independent detector indices},
project source, version 42,
\texttt{papers/pole\_parameters/}\allowbreak
\texttt{common\_parameter\_guard\_reduction\_20261116\_v42.tex}.
\bibitem{Local}
\emph{Local solubility at every shared pole parameter},
project source, version 41,
\texttt{papers/pole\_parameters/}\allowbreak
\texttt{shared\_parameter\_local\_solubility\_20261115\_v41.tex}.
\end{thebibliography}
\end{document}
